\family{Anacardiaceae}{Cashew family}
\label{fam:anacardiaceae}

Trees or shrubs with resin canals in the bark, so that the slash often
smells of turpentine or mango. Leaves alternate, simple or pinnate,
without stipules. Flowers small, usually 5-merous, in terminal or
axillary panicles. Fruit usually a drupe. Two species in the Arami
Uplands, one of them cultivated.

\section*{Key to the species}
\begin{florakey}
	\lead{1a}{Leaves simple, obovate, leathery; fruit a kidney-shaped nut
		on a swollen fleshy stalk}{\taxon{Anacardium occidentale}}
	\lead{1b}{Leaves pinnate with a terminal leaflet; fruit an ovoid
		yellow drupe}{\taxon{Spondias mombin}}
\end{florakey}

\species{Anacardium occidentale}{L.}
\label{sp:ANA01}

\field{Local and trade names} Yoruba: Kaju. Trade: Cashew.

\field{Status} Cultivated. Introduced from tropical America and now
widely planted in farms and plantations.

\field{Description} Evergreen tree to 12~m with a low, spreading
crown. Leaves simple, alternate, obovate, leathery, 8--20
$\times$ 6--12~cm, rounded or notched at the apex. Flowers small,
fragrant, pale with pink or red streaks, in terminal panicles. Fruit a
kidney-shaped nut 2--3~cm long, borne on a swollen, pear-shaped, red
or yellow fleshy stalk (the cashew apple).

\field{Phenology} Flowering December to March; fruiting February to
May.

\field{Habitat} Farms and plantations in the forest--savanna mosaic
and Guinea savanna woodland.

\field{Uses} The roasted kernel is eaten and sold, and the cashew
apple is eaten fresh or pressed for juice. The oil in the shell is
caustic and is used in industry \citep{burkill1985,orwa2009}.

\field{Specimens examined} Agbede (C2), 540~m, st., 7~Jul 1936, \specimen{Harding}{1534}{OKF};
Agbede (C2), 540~m, fr., 6~May 1937, \specimen{Harding}{1659}{OKF};
Agbede (C2), 540~m, fr., 19~Apr 1938, \specimen{Harding}{1960}{OKF};
Oke-Asa farmland (B4), 690~m, st., 5~Jun 1965, \specimen{Olaniyan}{745}{SVN};
Idofin (A3), 760~m, fr., 13~Apr 1967, \specimen{Olaniyan}{938}{SVN};
Asiri Crossing (A2), 740~m, fl., 6~Jan 2019, \specimen{Alade}{133}{AUH};
Asiri Crossing (A2), 740~m, fr., 8~Feb 2025, \specimen{Alade}{978}{AUH}.

\field{Distribution} Grid squares A2, A3, B4, C2.

\begin{center}
\distmap{A/2, A/3, B/4, C/2}
\end{center}

\species{Spondias mombin}{L.}
\label{sp:ANA02}

\field{Synonyms} \sciname{Spondias lutea}{L.}

\field{Local and trade names} Yoruba: Iyeye. Trade: Hog plum.

\field{Status} Native.

\field{Description} Deciduous tree to 25~m with thick, corky, deeply
fissured bark. Leaves imparipinnate with 5--10 pairs of leaflets;
leaflets with an unequal base and a long-pointed apex. Flowers small,
white or cream, in large terminal panicles. Fruit an ovoid drupe
3--4~cm long, yellow when ripe, with an acid, aromatic pulp.

\field{Phenology} Flowering February to April; fruiting June to
September.

\field{Habitat} Lowland semi-deciduous forest and forest--savanna
mosaic; common around villages, where it is planted from cuttings.

\field{Uses} The fruit is eaten fresh. Large cuttings root readily and
are used as living fence posts \citep{keay1989,burkill1985}.

\field{Specimens examined} Olowo Stream (D3), 440~m, fr., 7~Jul 1934, \specimen{Harding}{1173}{OKF};
Ireke Ferry (D4), 360~m, st., 15~Dec 1934, \specimen{Harding}{1260}{OKF};
Igbo-Oba Forest Reserve (C3), 470~m, fr., 15~Jun 1937, \specimen{Harding}{1710}{OKF};
Ewe Stream gallery forest (B3), 580~m, st., 27~Jan 1961, \specimen{Olaniyan}{262}{SVN};
Aramiwo Foot (B2), 640~m, fr., 19~Sep 1966, \specimen{Olaniyan}{853}{SVN};
Olowo Stream (D3), 440~m, st., 23~May 2022, \specimen{Alade}{611}{AUH};
Olowo Stream (D3), 440~m, fr., 9~Sep 2024, \specimen{Alade}{862}{AUH}.

\field{Distribution} Grid squares B2, B3, C3, D3, D4.

\begin{center}
\distmap{B/2, B/3, C/3, D/3, D/4}
\end{center}
